\chapter{結論}
\label{chap:conclusion}

本章は、モデル構成と実装上の到達点を整理する草稿である。実参加者による評価結果が未記載の段階であり、精度改善、意味分離、支援効果を実証した結論とはしない。

\section{本研究の総括}

研究用キーボードと任意同意の生活情報を用い、拠点、個人固有の行動パターン、現在・日次5尺度を区別する生活状態モデルを構成した。固定24活動分類と独立Task推定を重複させるのではなく、一つの三階層LSTMから各表現を得る。24活動の既存回答は削除せず、学習後の意味注釈と従来方式の比較へ役割を限定した。

観測、推定、本人確認、命名、訂正を別の記録として保持する。特徴が似たまとまりが得られたことと、そのまとまりの意味が分かったことを区別し、未確認候補は候補のまま扱う。

\section{数理モデルと実装の到達点}

Micro・Meso・Macroの時間表現を共有し、PlaceとBehaviorには独立した埋め込み射影、現在・日次尺度には別の回帰出力を用いる。自己教師ありの再構成・将来予測、本人確認による関係損失、尺度教師を同時に最適化する構成を定義した。

行動表現には閾値型のプロトタイプ生成と時間遷移を適用する。5尺度の予測値を行動の距離へ直接加えず、訓練側で作ったモデルとマップを推論時に固定する。再学習時に元観測を再埋め込みし、古い特徴空間の重心を新しいモデルへ混ぜない。

受信時刻による入力制限、分割を跨ぐ将来教師の除外、訓練側だけの正規化、旧checkpointの保持を実装上の条件とした。状態保持型TBPTT、オンライン重み更新、参加者間の部分プーリングを、この新構成で実装済みの機能とはしない。

\section{探索的評価の位置付け}

最大5名・各14日間の予備研究では、まず取得の実現可能性と回答負担を検討する。拠点・行動の本人確認IDとの一致、未知・抑制の割合、過分割・混合、5尺度の参加者別誤差を分けて評価する計画である。

自動テストは、固定分類の混入、情報漏洩、保存・推論の回帰を確認する。実生活における意味の分離や推定性能を証明するものではない。観測や回答が少なく評価条件へ到達しなかった場合も、その範囲を結果として記録する必要がある。

\section{今後の課題}

PlaceとBehaviorの共起を超えた分離、新たに生成した候補の意味的な妥当性、候補の安定性、少数ラベル下の尺度推定を実データで検証する。比較方式には同じ評価時点と境界を適用し、追加アブレーションと参加者単位の集計を実装する。モデル版ごとの候補・ID・本人訂正の対応も追跡する。

新しい状態IDに対応する歩数基準と運動・休憩提案の検証は残るため、現段階の新方式では状態別歩数と旧疲労推定による提案判定を停止する。入力訂正の変化と本人確認による補助的な休憩提案経路は残るが、質問予算等の条件を満たす場合だけ動作し、新しい生活状態モデルによる支援効果を示すものではない。生活状態の表現を得る段階と、本人に有益な支援へつなぐ段階を分け、同意、負担、提示条件を確認した上で後続研究へ進める。
